% SPDX-License-Identifier: Apache-2.0
% covers: AxiPins
\datasheet{AxiPins}{A host's AXI4 pins joined to the AXI link's channels.}

\dsfacts{\code{lib/parts/src/bus/axi\_pins.rs}}%
{\code{txhdl\_parts::bus::axi\_pins::AxiPins}}%
{\code{//docs:pcie}}

\subsection*{Function}

A core from elsewhere that is an AXI4 host has pins: a \code{valid} and
a \code{ready} per channel and a wire per field, named as AXI4 names
them. \code{AxiPins} joins those pins to the link's five channels, so it
goes where an \code{AxiHost} would, with an \code{AxiPer} or a router
behind it. Its inputs are the struct \code{AxiPinsIn}: the host's 24
pins and the \code{b} and \code{r} channels from the link. Its outputs
are \code{AxiPinsOut}: the \code{aw}, \code{ar} and \code{w} channels
to the link, and the eleven pins the host reads.

\code{axi\_pins::sim} has \code{PinHost}, a host on the pins written as
a simulation, and \code{pin\_host} and \code{pins}, which make the
wires.

\subsection*{Parameters}

\begin{center}\footnotesize
\begin{tabular}{@{}lp{0.7\textwidth}@{}}
\toprule
Parameter & Meaning \\
\midrule
\code{A} & The address width, in bits. The tables use 32. \\
\code{D} & The data width, in bits. The tables use 64. \\
\code{S} & The strobe width, \code{D / 8}. The tables use 8. \\
\code{I} & The identifier width, in bits. The tables use 4. \\
\bottomrule
\end{tabular}
\end{center}

The tables are generated for \code{AxiPins<32, 64, 8, 4>}, the widths of
the PCIe endpoint's master as \code{PcieBar} uses it.

\dstables{AxiPins}

\subsection*{Behaviour}

\begin{itemize}
\item It holds no state; every output is a function of the step's
inputs and the channels.
\item \code{awready}, \code{arready} and \code{wready} are the room on
\code{aw}, \code{ar} and \code{w}.
\item When \code{awvalid} is high and \code{aw} has room, an address
phase goes onto \code{aw} with the pins' \code{id}, \code{addr},
\code{len}, \code{size}, \code{lock}, \code{cache} and \code{prot}, the
burst type decoded from its two bits, and \code{qos} and \code{region}
zero. The same holds for \code{arvalid} and \code{ar}, and for
\code{wvalid} and \code{w}, with \code{data}, \code{strb} and
\code{last}.
\item \code{bvalid} is high while \code{b} holds a response, with
\code{bid} and \code{bresp} its fields, the response encoded in two
bits. It is taken in a step \code{bready} is high.
\item \code{rvalid}, \code{rid}, \code{rdata}, \code{rresp} and
\code{rlast} are the head of \code{r} the same way, taken in a step
\code{rready} is high.
\end{itemize}

\subsection*{Verification}

\begin{itemize}
\item \code{a\_host\_on\_pins\_reaches\_a\_peripheral\_through\_the\_link}
in \code{bus::axi\_pins}, in \code{//lib/parts:parts\_test}, writes three
words through the pins, an \code{AxiPer} and a simulated RAM, reads four
back with their identifier, and reads past the RAM's end, answered
\code{SlvErr}.
\item \code{ex\_axi\_pins} runs \code{AxiPins<32, 64, 8, 4>} the same
way, and the build simulates its netlist against that run under nvc and
Verilator: \code{//docs:axi\_pins\_sim\_axi\_pins\_tb\_test} and
\code{//docs:axi\_pins\_vsim\_test}.
\item \code{//pcie:bar\_test} drives it inside \code{PcieBar}.
\end{itemize}

\subsection*{Used by}

\code{PcieBar}, as its field \code{pins}, between the PCIe endpoint's
master and a peripheral tracker. The example \code{ex\_axi\_pins}.

\subsection*{Limits}

Its \code{ready} outputs follow the channels' room in the same step, so
a path runs from a channel's state through the unit to the host's pins
with no register. It has no \code{AxQOS}, \code{AxREGION} or user pins,
and drives \code{qos} and \code{region} as zero.
